\subsection{EIGENVALUES OF A COXETER TRANSFORMATION}

Since all Coxeter transformations are conjugate (no. 1, Prop. 1), they all have the same characteristic polynomial P(T). Let h be the Coxeter number of W. Then

\[
\mathrm{P} (\mathrm{T}) = \prod_ {j = 1} ^ {l} \left(\mathrm{T} - \exp \frac {2 i \pi m _ {j}}{h}\right),
\]

where \(m_{1}, m_{2}, \ldots, m_{l}\) are integers such that

\[
0 \leqslant m _ {1} \leqslant m _ {2} \leqslant \dots \leqslant m _ {l} <   h.
\]

DEFINITION 2. The integers \(m_1, m_2, \ldots, m_l\) are called the exponents of W.

Let C be a chamber determined by \(\mathfrak{H}\) , \(\mathrm{H}_1, \ldots, \mathrm{H}_l\) its walls, and put \(s_i = s_{\mathrm{H}_i}\) . Denote by \(e_i\) the unit vector orthogonal to \(\mathrm{H}_i\) and on the same side of \(\mathrm{H}_i\) as C. By Prop. 2 of the Appendix of Chap. IV, we can assume that the \(\mathrm{H}_i\) are numbered so that \(e_1, e_2, \ldots, e_r\) are pairwise orthogonal and \(e_{r+1}, e_{r+2}, \ldots, e_l\) are pairwise orthogonal. Then \(s' = s_1 s_2 \ldots s_r\) is the orthogonal symmetry with respect to the subspace

\[
\mathrm{V} ^ {\prime} = \mathrm{H} _ {1} \cap \mathrm{H} _ {2} \cap \dots \cap \mathrm{H} _ {r},
\]

\(s^{\prime \prime} = s_{r + 1}s_{r + 2}\dots s_{l}\) is the orthogonal symmetry with respect to the subspace

\[
\mathrm{V} ^ {\prime \prime} = \mathrm{H} _ {r + 1} \cap \mathrm{H} _ {r + 2} \cap \dots \cap \mathrm{H} _ {l},
\]

and \(c = s's''\) is a Coxeter transformation. Since \((e_{1}, \ldots, e_{l})\) is a basis of V, V is the direct sum of \(V'\) and \(V''\) .

We deduce first that 1 is not an eigenvalue of c.~For if \(x \in V\) is such that \(c(x) = x\) , then \(s'(x) = s''(x)\) , so \(x - s'(x) = x - s''(x)\) is orthogonal to \(V'\) and \(V''\) , and hence is zero. Thus, \(x = s'(x) = s''(x) \in V' \cap V'' = \{0\}\) .

Consequently,

\[
0 <   m _ {1} \leqslant m _ {2} \leqslant \dots \leqslant m _ {l} <   h.\tag{1}
\]

The characteristic polynomial of c has real coefficients. Thus, for all j, the power of \(T - \exp\frac{2i\pi m_{j}}{h}\) in \(P(T)\) is equal to that of \(T - \exp\frac{2i\pi(h - m_{j})}{h}\) . Hence

\[
m _ {j} + m _ {l + 1 - j} = h \quad (1 \leqslant j \leqslant l).\tag{2}
\]

Adding the equalities (2) we obtain

\[
m _ {1} + m _ {2} + \dots + m _ {l} = \frac {1}{2} l h.\tag{3}
\]

Lemma 2. Assume that W is irreducible and that \(\dim V \geqslant 2\) . With the preceding notation, there exist two linearly independent vectors \(z', z''\) such that

\begin{enumerate}
\def\labelenumi{(\roman{enumi})}
\item
  the plane P generated by \(z', z''\) is stable under \(s'\) and \(s''\) ;
\item
  \(s'|\mathrm{P}\) and \(s''|\mathrm{P}\) are orthogonal reflections with respect to \(\mathbf{R}z'\) and \(\mathbf{R}z''\) .
\item
  \(z', z'' \in \overline{\mathbf{C}}\) , and \(P \cap C\) is the set of linear combinations of \(z', z''\) with coefficients \(> 0\) .
\end{enumerate}

Let \((e^{1},\ldots,e^{l})\) be the basis of V such that \((e^{i}|e_{j})=\delta_{ij}\) . Then C is the open simplicial cone determined by the \(e^{i}\) ( \(\S 3\) , no. 9, Prop. 7). It is clear that \(V'\) is generated by \(e^{r+1},\ldots,e^{l}\) and \(V''\) by \(e^{1},\ldots,e^{r}\) . Let q be the endomorphism of V such that \(q(e^{1})=e_{1},\ldots,q(e^{l})=e_{l}\) . Its matrix with respect to \((e^{1},\ldots,e^{l})\) is \(Q=((e_{i}|e_{j}))\) . We have \((e_{i}|e_{j})\leqslant0\) for \(i\neq j\) ( \(\S 3\) , no. 4, Prop. 3). Since W is irreducible, there does not exist any partition \(\{1,2,\ldots,l\}=I_{1}\cup I_{2}\) such that \((e_{i}|e_{j})=0\) for \(i\in I_{1}\) and \(j\in I_{2}\) . Thus ( \(\S 3\) , no. 5, Lemma 4) Q has an eigenvector \((a_{1},\ldots,a_{l})\) all of whose coordinates are \textgreater0; let a be the corresponding eigenvalue. Put

\[
\begin{array}{r} z = a _ {1} e ^ {1} + \dots + a _ {l} e ^ {l}, \\ z ^ {\prime \prime} = a _ {1} e ^ {1} + \dots + a _ {r} e ^ {r} \in V ^ {\prime \prime} \cap \overline {{C}}, \\ z ^ {\prime} = a _ {r + 1} e ^ {r + 1} + \dots + a _ {l} e ^ {l} \in V ^ {\prime} \cap \overline {{C}}, \end{array}
\]

and let P be the plane generated by \(z'\) and \(z''\) . Then \(P \cap C\) is the set of linear combinations of \(z'\) and \(z''\) with coefficients \(> 0\) . The relation \(q(z) = az\) gives \(\sum_{j=1}^{l} a_j e_j = \sum_{j=1}^{l} aa_j e^j\) ; scalar multiplying by \(e_k\) (where \(k \leqslant r\) ) gives \(a_k + \sum_{j=r+1}^{l} a_j (e_j | e_k) = aa_k\) ; thus

\[
\begin{array}{l} (a - 1) z ^ {\prime \prime} = \sum_ {k = 1} ^ {r} \left(\sum_ {j = r + 1} ^ {l} a _ {j} (e _ {j} | e _ {k})\right) e ^ {k} \\ \qquad = \sum_ {j = r + 1} ^ {l} a _ {j} \left(\sum_ {k = 1} ^ {r} (e _ {j} | e _ {k}) e ^ {k}\right) \\ \qquad = \sum_ {j = r + 1} ^ {l} a _ {j} \left(- e ^ {j} + \sum_ {k = 1} ^ {l} (e _ {j} | e _ {k}) e ^ {k}\right) \\ \qquad = - \sum_ {j = r + 1} ^ {l} a _ {j} e ^ {j} + \sum_ {j = r + 1} ^ {l} a _ {j} e _ {j} \\ \qquad = - z ^ {\prime} + \sum_ {j = r + 1} ^ {l} a _ {j} e _ {j}. \end{array}
\]

Thus, \((a-1)z'' + z'\) is orthogonal to \(e^{1}, \ldots, e^{r}\) , that is, to \(V''\) . Hence, \(s''\) leaves stable the plane generated by \(z''\) and \((a-1)z'' + z'\) , that is, P. Similarly, \(s'\) leaves P stable. Since \(z' \in P \cap V'\) and \(z'' \in P \cap V''\) , \(s'|P\) and \(s''|P\) are the reflections with respect to \(Rz'\) and \(Rz''\) .

THEOREM 1. Assume that W is irreducible. Then:

\begin{enumerate}
\def\labelenumi{(\roman{enumi})}
\item
  \(m_{1} = 1, m_{l} = h - 1\) .
\item
  \(\operatorname{Card}(\mathfrak{H}) = \frac{1}{2} lh\) .
\end{enumerate}

We retain the preceding notation. The restriction of \(c = s's''\) to \(P\) is the rotation with angle \(2(z'', z')\) (§2, no. 5, Cor. of Prop. 6). Since \(c\) is of order \(h\) , the \(h\) elements \(1, c, \ldots, c^{h-1}\) of \(W\) are pairwise distinct; the elements \(s', s'c, \ldots, s'c^{h-1}\) are thus pairwise distinct, and are distinct from the preceding elements since \(c^i |P\) is a rotation and \(s'c^j |P\) is a reflection. The set

\[
\{1, c, \dots , c ^ {h - 1}, s ^ {\prime}, s ^ {\prime} c, \dots , s ^ {\prime} c ^ {h - 1} \}
\]

is the subgroup \(W'\) of W generated by \(s'\) and \(s''\) , and induces on P the group \(W''\) generated by the orthogonal reflections with respect to \(Rz', Rz''\) . The transform of C by an element of \(W'\) is either disjoint from -C or equal to -C. Thus, the transform of \(P \cap C\) by an element of \(W''\) is either disjoint from \(-(P \cap C)\) or equal to \(-(P \cap C)\) . Hence, for a suitable orientation of \(P\) , there exists an integer \(m > 0\) such that \((z'', z') = \frac{\pi}{m}\) (§2, no. 5, Cor. of Prop. 7). Moreover, the sets \(g'(C)\) , for \(g' \in W'\) , are pairwise disjoint; the sets \(g''(P \cap C)\) , for \(g'' \in W''\) , are thus pairwise disjoint; so \(W''\) is of order \(2h\) . Hence \(m = h\) . By definition, \(c|P\) is a rotation with angle \(\frac{2\pi}{m}\) , and so has eigenvalues \(\exp \frac{2i\pi}{h}, \exp \frac{2i\pi(h - 1)}{h}\) . This proves that \(m_1 = 1, m_l = h - 1\) .

The transforms of \(\mathbf{R}z'\) and \(\mathbf{R}z''\) by \(W'\) are \(h\) lines \(D_1, \ldots, D_h\) of \(P\) , and the points of \(P - (D_1 \cup \cdots \cup D_h)\) are transforms by the elements of \(W'\) of points of \(P \cap C\) . Thus, a hyperplane of \(\mathfrak{H}\) necessarily cuts \(P\) along one of the lines \(D_i\) , and consequently is a transform by an operation of \(W'\) of a hyperplane of \(\mathfrak{H}\) containing \(\mathbf{R}z'\) or \(\mathbf{R}z''\) .

Now, any \(H \in \mathfrak{H}\) that contains \(Rz'\) is one of the hyperplanes \(H_{1}, \ldots, H_{r}\) . Indeed, let \(e_{H}\) be the unit vector orthogonal to H and on the same side of H as C. Then \(e_{H} = \lambda_{1}e_{1} + \cdots + \lambda_{l}e_{l}\) with the \(\lambda_{i}\) all \(\geqslant 0\) (§3, no. 5, Lemma 6, (i)). Now \(0 = (e_{\mathrm{H}}|z') = \lambda_{r+1}a_{r+1} + \cdots + \lambda_{l}a_{l}\) , so

\[
\lambda_ {r + 1} = \dots = \lambda_ {l} = 0, \quad \text { and } \quad e _ {\mathrm{H}} = \lambda_ {1} e _ {1} + \dots + \lambda_ {r} e _ {r}.
\]

Suppose that two of the \(\lambda_{i}\) were non-zero, for example \(\lambda_{1}\) and \(\lambda_{2}\) ; since \(e_{1},\ldots,e_{r}\) are pairwise orthogonal, we would have

\[
s _ {1} (e _ {\mathrm{H}}) = - \lambda_ {1} e _ {1} + \lambda_ {2} e _ {2} + \dots + \lambda_ {r} e _ {r}
\]

and the coordinates of \(s_{1}(e_{\mathrm{H}})\) would not be of the same sign, which is absurd (loc. cit.). Thus, \(e_{H}\) is proportional to one of the vectors \(e_{1},\ldots,e_{r}\) , which proves our assertion. Similarly, any \(H\in \mathfrak{H}\) that contains \(Rz''\) is one of the hyperplanes \(H_{r+1},\ldots,H_{l}\) .

The number of elements of \(\mathfrak{H}\) containing \(\mathbf{R}z'\) or \(\mathbf{R}z''\) is therefore \(l\) . If \(h\) is even, \(\operatorname{Card}(\mathfrak{H})\) is thus equal to \(\frac{h}{2} l\) . If \(h\) is odd, \(\operatorname{Card}(\mathfrak{H})\) is equal to \(\frac{h - 1}{2} l + r\) , and also to \(\frac{h - 1}{2} l + (l - r)\) ; hence \(r = l - r\) , so \(r = \frac{l}{2}\) , and \(\operatorname{Card}(\mathfrak{H}) = \frac{h - 1}{2} l + \frac{l}{2} = \frac{h}{2} l\) .

Remark. Retain the notation of the preceding proof. Let \(c'\) be the \(\mathbf{C}\) -linear extension of \(c\) to \(V \otimes_{\mathbf{R}} \mathbf{C}\) , and \(c''\) the restriction of \(c'\) to \(P \otimes_{\mathbf{R}} \mathbf{C}\) . From our study of \(c|P\) , \(c''\) has an eigenvector \(x\) corresponding to the eigenvalue \(\exp \frac{2i\pi}{h}\) , and this eigenvector does not belong to any of the sets \(D \otimes_{\mathbf{R}} \mathbf{C}\) , where \(D\) denotes a line of \(P\) (since \(D\) is not stable under \(c\) ). Now, for any \(H \in \mathfrak{H}\) , we have seen that \(H \cap P\) is a line; thus, \(x \notin H \otimes_{\mathbf{R}} \mathbf{C}\) .

COROLLARY. Let \(R_{0}\) be the set of unit vectors of V orthogonal to an element of H. If W is irreducible,

\[
\sum_ {u \in R _ {0}} (x | u) ^ {2} = h (x | x)\tag{4}
\]

for all \(x \in V\) .

Put \(f(x) = \sum_{u \in R_0} (x|u)^2\) . It is clear that \(f\) is a positive quadratic form invariant under \(W\) , and non-degenerate since the \(e_i\) form a basis of \(V\) . Since

W is irreducible, there exists a constant \(\beta\) such that \(f(x) = \beta (x|x)\) (§ 2, no. 1, Prop. 1). If \((x_{i})_{1\leqslant i\leqslant l}\) is an orthonormal basis of V for the scalar product \((x|y)\) , then

\[
\begin{array}{l} \beta l = \sum_ {i = 1} ^ {l} \beta (x _ {i} | x _ {i}) = \sum_ {i = 1} ^ {l} f (x _ {i}) = \sum_ {i = 1} ^ {l} \sum_ {u \in R _ {0}} (x _ {i} | u) ^ {2} \\ = \sum_ {u \in R _ {0}} 1 = \operatorname{Card} (R _ {0}) = 2 \operatorname{Card} (\mathfrak {H}) = h l. \end{array}
\]

Thus \(\beta = h\) , which proves (4).

PROPOSITION 2. If \(W\) is irreducible and \(h\) is even, the unique element of \(W\) that transforms \(C\) to \(-C\) is \(c^{h/2}\) .

We use the notation in the proof of Th. 1. Since \(c|P\) is a rotation through an angle \(\frac{2\pi}{h}\) , \(c^{h/2}\) transforms \(z'\) to \(-z'\) , \(z''\) to \(-z''\) , and hence \(z' + z'' = z\) to \(-z\) . Now \(z \in C\) , so the chamber \(c^{h/2}(C)\) is necessarily \(-C\) .

PROPOSITION 3. Assume that W is irreducible. Let \(u_1, \ldots, u_l\) be homogeneous elements of the symmetric algebra \(S = \mathbf{S}(V)\) , algebraically independent over \(\mathbf{R}\) and generating the algebra of elements of \(S\) invariant under \(W\) (§5, no. 3, Th. 3). If \(p_j\) is the degree of \(u_j\) , the exponents of \(W\) are \(p_1 - 1, \ldots, p_l - 1\) .

Put \(V' = V \otimes_{\mathbf{R}} \mathbf{C}\) , \(S' = \mathbf{S}(V') = S \otimes_{\mathbf{R}} \mathbf{C}\) , and extend the scalar product on \(V\) to a hermitian form on \(V'\) . If \(c\) is a Coxeter transformation of \(W\) , there exists an orthonormal basis \((X_i)_{1 \leqslant i \leqslant l}\) of \(V'\) consisting of eigenvectors of \(c \otimes 1\) (Algebra, Chap. IX, §7, no. 3, Prop. 4); moreover, we can assume that, for \(1 \leqslant j \leqslant l\) , \(X_j\) corresponds to the eigenvalue \(\exp \frac{2i\pi m_j}{h}\) of \(c \otimes 1\) . It is clear that \(S'\) can be identified with the algebra \(\mathbf{C}[X_1, \ldots, X_l]\) , and that we can write \(u_j \otimes 1 = f_j(X_1, \ldots, X_l)\) , where \(f_j\) is a homogeneous polynomial of degree \(p_j\) in \(\mathbf{C}[X_1, \ldots, X_l]\) . Put \(D_j = \frac{\partial}{\partial X_j}\) and \(J = \det(D_k f_j)\) . Recall (§5, no. 4, Prop. 5) that \(J(X_1, \ldots, X_l)\) is proportional to the product in \(S'\) of \(\text{Card}(\mathfrak{H})\) vectors \(y_k\) of \(V\) each of which is orthogonal to a hyperplane of \(\mathfrak{H}\) . Since we can assume that \(X_1 \notin H \otimes C\) for all \(H \in \mathfrak{H}\) (Remark), the \(X_1\) component of each of the vectors \(y_k\) is non-zero, so \(J(1, 0, \ldots, 0) \neq 0\) . The rule for expanding determinants now proves the existence of a permutation \(\sigma\) of \(\{1, 2, \ldots, l\}\) such that \((D_{\sigma(j)} f_j)(1, 0, \ldots, 0) \neq 0\) for all \(j\) . Since \(D_{\sigma(j)} f_j\) is homogeneous of degree \(p_j - 1\) , the coefficient of \(X_1^{p_j - 1} X_{\sigma(j)}\) in \(f_j(X_1, \ldots, X_l)\) is non-zero. Now \(f_j(X_1, \ldots, X_l)\) is invariant under \(c \otimes 1\) , and

\[
(c \otimes 1) (\mathrm{X} _ {1} ^ {p _ {j} - 1} \mathrm{X} _ {\sigma (j)}) = \left(\exp \frac {2 i \pi}{h} (p _ {j} - 1 + m _ {\sigma (j)})\right) (\mathrm{X} _ {1} ^ {p _ {j} - 1} \mathrm{X} _ {\sigma (j)}).
\]

This proves that \(p_j - 1 + m_{\sigma(j)} \equiv 0 \pmod{h}\) . Now \(h - m_{\sigma(j)}\) is an exponent (formula (2)). Permuting the \(u_j\) if necessary, we can assume that \(p_j - 1 \equiv m_j \pmod{h}\) for all \(j\) . Since \(p_j - 1 \geqslant 0\) and \(m_j < h\) , we have \(p_j - 1 = m_j + \mu_j h\) with \(\mu_j\) an integer \(\geqslant 0\) . By §5, Prop. 3, we see that

\[
\operatorname{Card} (\mathfrak {H}) = \sum_ {j = 1} ^ {l} (p _ {j} - 1) = \sum_ {j = 1} ^ {l} m _ {j} + h \sum_ {j = 1} ^ {l} \mu_ {j}.
\]

Taking into account formula (3) and Th. 1 (ii), we obtain \(h \sum_{j=1}^{l} \mu_j = 0\) , so \(\mu_j = 0\) for all \(j\) , and finally \(p_j - 1 = m_j\) for all \(j\) .

COROLLARY 1. If \((m_i)_{1\leqslant i\leqslant l}\) is the increasing sequence of exponents of W, the order of W is equal to \((m_1 + 1)(m_2 + 1)\ldots (m_l + 1)\) .

This follows from the relations \(m_j + 1 = p_j\) and §5, Cor. of Th. 3.

COROLLARY 2. If \(c\) is a Coxeter transformation of \(W\) ,

\[
\exp \left(\frac {2 i \pi}{h}\right) \quad \text { and } \quad \exp \left(\frac {- 2 i \pi}{h}\right)
\]

are eigenvalues of \(c\) of multiplicity 1.

Otherwise, there would exist two non-proportional homogeneous invariants of degree 2 in S, and hence two non-proportional quadratic forms on V* invariant under W, contrary to §2, no. 1, Prop. 1.

COROLLARY 3. The homothety with ratio -1 of V belongs to W if and only if all the exponents of W are odd. In that case, h is even and \(c^{h/2} = -1\) for any Coxeter transformation c of W.

The first assertion follows from §5, no. 3, Prop. 4. Assume that the exponents of \(W\) are odd. Then \(h\) is even by formula (2), and

\[
\left(\exp \frac {2 i \pi m _ {j}}{h}\right) ^ {h / 2} = \exp (i \pi m _ {j}) = - 1;
\]

thus \(c^{h/2} = -1\) since \(c\) is a semi-simple automorphism of \(V\) (Algebra, Chap. IX, §7, no. 3, Prop. 4).

